% =============================================================================
% SECTION 2: BACKGROUND AND RELATED WORK
% =============================================================================
\section{Background and Related Work}

\subsection{Shellcode and Loader: a What/How Decomposition}

Modern malware that delivers arbitrary code commonly separates two concerns, and we adopt this separation throughout the paper.
The \textit{shellcode} answers \textbf{what} the malware does: it is a self-contained, position-independent code fragment whose functionality may be, for instance, establishing a reverse shell, encrypting files, self-replicating, or exfiltrating data.
The \textit{loader} answers \textbf{how} that shellcode reaches execution: it is the code that delivers, transforms, places, and runs the shellcode inside a victim environment while attempting to bypass defenses.
Under this lens, the shellcode is essentially interchangeable---the same reverse-TCP shellcode can be paired with many different loaders, and a single loader can deliver many different shellcodes.
This interchangeability is why modern offensive toolchains and defensive analyses alike treat loader engineering and payload engineering as distinct problems.

The present work focuses on the \textit{how}.
Our platform and experiments are designed around the shellcode/loader separation and study the loader as the primary unit of variation, which in turn motivates the stage-based decomposition introduced in Section~3.

\subsection{Open-Source Loader Frameworks}

Several open-source frameworks generate shellcode loaders for red-team operations and constitute the closest practical neighbors of this work.
Donut~\cite{donut} converts .NET assemblies or native PE files into position-independent shellcode packaged with a loader stub, and is widely reused as a building block.
ScareCrow~\cite{scarecrow} focuses on EDR bypass through direct syscalls, DLL unhooking, and sideloaded entry points.
Inceptor~\cite{inceptor} provides a template-driven build system for AV/EDR evasion, supporting multiple encoders, encryption primitives, and injection techniques.
Havoc~\cite{havoc} is a post-exploitation command-and-control framework whose implant includes a configurable loader module with encoding and syscall options.

These tools are valuable for operational use but are not organized for systematic research.
They typically expose a flat option surface rather than a per-stage decomposition, bundle multiple concerns inside opaque templates, and are not paired with a controlled experimental harness that isolates the effect of each internal choice on detection outcomes.
As a result, they provide neither a vocabulary nor a reproducible test fixture suitable for the kind of stage-level comparative study we aim to enable.

\subsection{Individual Loader Techniques in the Literature}

Research on individual loader techniques is extensive but fragmented.
Hosseini~\cite{process_injection} surveys ten process injection techniques and discusses their detection surfaces.
Chernenko~\cite{api_unhooking} analyzes user-mode hook bypass and direct invocation of system calls as red-team tradecraft.
The \textit{Hell's Gate} technique~\cite{hellsgate} demonstrates how to dynamically resolve system service numbers from a loaded \texttt{ntdll} copy, enabling direct-syscall loaders that evade user-mode hooks.
SysWhispers~\cite{syswhispers} and its successors generate syscall stubs for large subsets of the Windows NT API to support the same goal.

These works each focus on one technique or technique family in isolation.
They do not place their contribution inside a broader loader pipeline, do not provide a framework that relates several techniques across stages, and do not offer an automated way to generate and compare technique variants.

\subsection{Malware Testing and Detection Platforms}

Automated malware analysis platforms such as Cuckoo Sandbox~\cite{cuckoo} and ANY.RUN~\cite{anyrun} execute submitted binaries in instrumented environments and return behavioral reports.
They are oriented toward analyzing pre-existing samples rather than \emph{generating} technique-level variants for controlled comparison, and they treat each sample as an opaque unit rather than as a composition of known stages.

Consumer-facing evaluations by AV-TEST~\cite{avtest} and AV-Comparatives~\cite{avcomparatives} measure detection rates of security products across malware corpora.
These evaluations emphasize aggregate detection scores and are not designed to attribute detection outcomes to specific internal techniques of a loader.

MITRE ATT\&CK~\cite{mitre_attack} supplies a widely adopted vocabulary for adversary tactics and techniques (e.g., T1055 Process Injection, T1620 Reflective Code Loading).
ATT\&CK operates at a strategic, cross-sample level and does not model the internal execution pipeline of a single loader binary; our pipeline model is intended to be complementary to ATT\&CK rather than a replacement.

\subsection{Research Gap}

Table~\ref{tab:gap} positions this work against representative systems in the three categories above.
We retain only the axes that are meaningful for a loader-focused study: stage-level decomposition of loader internals, automated generation of technique variants, and a stage-aware detection methodology.

\begin{table}[htbp]
    \centering
    \caption{Positioning against representative systems. Only this work combines stage-level decomposition, modular variant generation, and a stage-aware detection methodology; the last item is \textit{proposed} here and scoped to future work for empirical validation.}
    \label{tab:gap}
    \small
    \begin{tabular}{@{}p{2.8cm}ccc@{}}
        \toprule
        \textbf{System} & \rotatebox{70}{\textbf{Stage decomp.}} & \rotatebox{70}{\textbf{Variant gen.}} & \rotatebox{70}{\textbf{Stage-aware det.}} \\
        \midrule
        MITRE ATT\&CK~\cite{mitre_attack}         & Partial    & ---        & ---        \\
        Cuckoo / ANY.RUN~\cite{cuckoo,anyrun}     & ---        & ---        & ---        \\
        AV-TEST / AV-Comp.~\cite{avtest,avcomparatives} & ---  & ---        & ---        \\
        Donut~\cite{donut}                        & ---        & \checkmark & ---        \\
        ScareCrow~\cite{scarecrow}                & ---        & \checkmark & ---        \\
        Inceptor~\cite{inceptor}                  & ---        & \checkmark & ---        \\
        Havoc loader~\cite{havoc}                 & ---        & \checkmark & ---        \\
        \textbf{This work}                        & \checkmark & \checkmark & Proposed   \\
        \bottomrule
    \end{tabular}
\end{table}

Existing loader frameworks generate functional loaders but do not expose an internal stage structure for research.
Individual-technique studies explain specific mechanisms but are not framed within a unified pipeline.
Testing platforms consume pre-built samples and return aggregate verdicts rather than technique-level comparisons.
To our knowledge, no prior work combines a stage-based decomposition of shellcode loader internals with automated variant generation and a stage-aware detection methodology intended to make per-stage detection surfaces explicit.
